% from ega4-12-fr.tex
\oldpage[IV-3]{187}
exacte (cette assertion découlant de l'hypothèse pour $k=0$). Soit donc
$x\in Q_k$ un élément tel que $v_k(x)=0$. Comme par l'hypothèse de récurrence
la suite
$P_{k-1}\xrightarrow{u_{k-1}}Q_{k-1}\xrightarrow{v_{k-1}}R_{k-1}$ est exacte,
l'image canonique de $x$ dans
$Q_{k-1}=Q_k/(\mathfrak J^kQ/\mathfrak J^{k+1}Q)$ appartient à
$\operatorname{Im}(u_{k-1})$, donc il existe $z\in P_k$ tel que
$x=u_k(z)+x'$ avec $x'\in\mathfrak J^kQ/\mathfrak J^{k+1}Q$. Comme on a
$v_k\circ u_k=0$, la relation $v_k(x)=0$ entraîne $v_k(x')=0$, et par ailleurs
on a évidemment $v_k(x')\in\mathfrak J^kR/\mathfrak J^{k+1}R$~; tout revient
donc à prouver que la suite
\[
  \mathfrak J^kP/\mathfrak J^{k+1}P\xrightarrow{u'}
  \mathfrak J^kQ/\mathfrak J^{k+1}Q\xrightarrow{v'}
  \mathfrak J^kR/\mathfrak J^{k+1}R
\]
(où $u'$ et $v'$ proviennent de $u$ et $v$ par restriction et passage aux
quotients) est exacte. Or, par hypothèse,
$\mathfrak J^k/\mathfrak J^{k+1}$ est un $(A/\mathfrak J)$-module plat, donc la
suite
\[
  (P/\mathfrak JP)\otimes_{A/\mathfrak J}
  (\mathfrak J^k/\mathfrak J^{k+1})\xrightarrow{u''}
  (Q/\mathfrak JQ)\otimes_{A/\mathfrak J}
  (\mathfrak J^k/\mathfrak J^{k+1})\xrightarrow{v''}
  (R/\mathfrak JR)\otimes_{A/\mathfrak J}
  (\mathfrak J^k/\mathfrak J^{k+1})
\]
est exacte. Mais
$(Q/\mathfrak JQ)\otimes_{A/\mathfrak J}
(\mathfrak J^k/\mathfrak J^{k+1})
=Q\otimes_A(\mathfrak J^k/\mathfrak J^{k+1})$ s'identifie à
$\mathfrak J^kQ/\mathfrak J^{k+1}Q$ en vertu de la platitude de $Q$ sur $A$,
et de même
$(R/\mathfrak JR)\otimes_{A/\mathfrak J}
(\mathfrak J^k/\mathfrak J^{k+1})$ s'identifie à
$\mathfrak J^kR/\mathfrak J^{k+1}R$. Enfin, l'image de $u'$ s'identifie à
celle de $u''$, et cela termine la démonstration.
\end{proof}

\begin{corollary}[12.3.4]
\label{IV.12.3.4-fr}
Soient $f:X\to S$ un morphisme plat localement de présentation finie,
$\mathcal F$, $\mathcal G$ deux $\mathcal O_X$-Modules de présentation finie et
$f$-plats. Soient $n$ un entier $\geq0$, $U$ l'ensemble des $x\in X$ tels que
\[
  \bigl(\mathscr{E}\!xt^n_{\mathcal O_{X_{f(x)}}}
  (\mathcal F_{f(x)},\mathcal G_{f(x)})\bigr)_x=0
\]
(resp.
$\bigl(\mathscr{T}\!or_n^{\mathcal O_{X_{f(x)}}}
(\mathcal F_{f(x)},\mathcal G_{f(x)})\bigr)_x=0$). Alors $U$ est ouvert, et
l'on a
\[
  \mathscr{E}\!xt^n_{\mathcal O_X}(\mathcal F,\mathcal G)|U=0
\]
(resp.
$\mathscr{T}\!or_n^{\mathcal O_X}(\mathcal F,\mathcal G)|U=0$).
\end{corollary}

\begin{proof}
On peut évidemment se borner au cas où $S$ et $X$ sont affines. Alors (12.3.2),
il existe une résolution gauche $\mathcal L_\bullet=(\mathcal L_i)$ de
$\mathcal F$ formée de $\mathcal O_X$-Modules \emph{libres de type fini}. Si
l'on pose $\mathcal M^\bullet=\mathscr{H}\!om(\mathcal L_\bullet,\mathcal G)$
(resp. $\mathcal N_\bullet=\mathcal L_\bullet\otimes_{\mathcal O_X}\mathcal G$),
on en déduit que chacun des $\mathcal M^i$ (resp. $\mathcal N_i$) est isomorphe
à un $\mathcal O_X$-Module de la forme $\mathcal G^m$, donc les $\mathcal M^i$
et $\mathcal N_i$ sont des $\mathcal O_X$-Modules de présentation finie et
$f$-plats. En outre, pour tout changement de base $S'\to S$, si l'on pose
$X'=X\times_SS'$, $\mathcal F'=\mathcal F\otimes_SS'$,
$\mathcal G'=\mathcal G\otimes_SS'$, $\mathcal L'_\bullet
=\mathcal L_\bullet\otimes_SS'$, $\mathcal L'_\bullet$ est encore une
résolution gauche de $\mathcal F'$ par des $\mathcal O_{X'}$-Modules libres de
type fini (2.1.10), et $\mathcal M'^\bullet=\mathcal M^\bullet\otimes_SS'$ est
égal à $\mathscr{H}\!om(\mathcal L'_\bullet,\mathcal G')$ (resp.
$\mathcal N'_\bullet=\mathcal N_\bullet\otimes_SS'$ est égal à
$\mathcal L'_\bullet\otimes_{\mathcal O_{X'}}\mathcal G'$), d'après ce qui
précède. En particulier, on a, pour tout $s\in S$,
\[
  \mathscr{E}\!xt^n_{\mathcal O_{X_s}}(\mathcal F_s,\mathcal G_s)
  =\mathcal H^n((\mathcal M^\bullet)_s)
\]
(resp.
$\mathscr{T}\!or_n^{\mathcal O_{X_s}}(\mathcal F_s,\mathcal G_s)
=\mathcal H_n((\mathcal N_\bullet)_s)$). Appliquant (12.3.3) aux complexes de
Modules $f$-plats $\mathcal M^\bullet$ et $\mathcal N_\bullet$, on en déduit
aussitôt le corollaire.
\end{proof}

% IV3_B26_P187_END_BEFORE_SECTION_13

% (continued in ega4-13-fr.tex, before its first page marker)
\section{Morphismes équidimensionnels}
\label{section:IV.13-fr}

Ce paragraphe est consacré à l'étude de la variation de la \emph{dimension}
des fibres d'un morphisme localement de type fini $f:X\to Y$ (qui est déjà
intervenue à propos de la «~formule des dimensions~» dans 5.5 et 5.6). On
prouve d'abord (13.1.3) que la fonction
$x\mapsto\dim_x(f^{-1}(f(x)))$ est toujours \emph{semi-continue
supérieurement} dans $X$ (théorème de semi-continuité de Chevalley). On étudie
ensuite plus particulièrement les morphismes, dits «~équidimensionnels~», pour
lesquels cette fonction est \emph{localement constante}.

